\documentclass[12pt]{book}

\usepackage{amsmath}
\usepackage{amsthm}
\usepackage{tikz}
\usepackage{caption}
\usepackage{subcaption}
\usepackage[hidelinks]{hyperref}
\usepackage{cleveref}
\usepackage{amssymb}


\newcommand{\Price}[2]{P(#1,#2)}
\newcommand{\Ptotal}[1]{P(#1)}


\newtheorem{theorem}{Theorem}
\newtheorem{proposition}{Proposition}
\newtheorem{definition}{Definition}
\newtheorem{corollary}{Corollary}
\newtheorem{lemma}{Lemma}
\newtheorem{ex}{Example}
\newtheorem*{pf}{Proof}

\begin{document}

\section{Convex functions in general}

\begin{definition}
    We say that a function $f$ defined over a domain $D$ of a euclidean space
        is \emph{convex} if for every two points $a,b\in D$ and every $0<
        \alpha<1$ with $\alpha a + (1-\alpha) b \in D$ it satisfies the
        inequality 
        \begin{equation*}
            f(\alpha a + (1-\alpha) b)
            \le
            \alpha f(a) + (1-\alpha) f(b)\,.
        \end{equation*}
    Moreover the function is said to be \emph{strictly} so if the inequality
        is never tight.
\end{definition}

\begin{lemma}
    Sum of two (strictly) convex function is (strictly) convex.
\end{lemma}
\begin{pf}
    Direct consequence of the definition.
    \qed
\end{pf}

\begin{lemma}
    Let $f_{1},f_{2}\colon D\to\mathbb{R}$ be two (strictly) convex function
        over a common euclidean domain $D$.
    Their maximum, which is the function
        \begin{equation*}
            f_{1}\vee f_{2}
            \colon
            D\to\mathbb{R}
            \colon
            x\mapsto\max\{f_{1}(x),f_{2}(x)\}
        \end{equation*}
        is (strictly) convex as well.
\end{lemma}
\begin{pf}
    For the sake of simplicity, within this proof, let us write $f$ in place
        of $f_{1}\vee f_{2}$.
    Let us consider $a,b\in D$ and $0<\alpha<1$ as in the definition of
        convexity.
    From the construction of $f$ it is clear that $f(\alpha a + (1-\alpha) b)$
        is either $f_{1}(\alpha a + (1-\alpha) b)$ or
        $f_{2}(\alpha a + (1-\alpha) b)$.
    We can assume, without any loss of generality, the the former is the case
        (were it the other way round, the proof can be carried out analogously
        just by replacing $f_{1}$ for $f_{2}$).
    Since, by assumption, $f_{1}$ is convex, we have 
        \begin{equation*}
            f_{1}(\alpha a + (1-\alpha) b)
            \le
            \alpha f_{1}(a) + (1-\alpha) f_{1}(b)\,.
        \end{equation*}
    From the construction of $f$ it is also clear that $f>f_{1}$ which, when
        combined with the latter inequality, yields
        \begin{equation*}
            f_{1}(\alpha a + (1-\alpha) b)
            \le
            \alpha f(a) + (1-\alpha) f(b)\,.
        \end{equation*}
    Finally, the convexity of $f$ is concluded just byt remembering that the
        left hand side is exactly $f(\alpha a + (1-\alpha) b)$.

    Now, if the convexity of both functions happens to be trict, the convexity
        inequality at the beginning of the consideratin becomes sharp and the
        sharpness is preserved in the following steps which proves the strict
        variant of the lemma.
    \qed
\end{pf}

\begin{lemma}
    Level cuts of (strictly) convex functions defined over convex domains are
        (strictly) convex sets.
\end{lemma}
\begin{pf}
    Let us omitt this for now.
    \qed
\end{pf}

As an example, let us pick an arbitratry point $c$ from the euclidean plane and
    consider the function
    \begin{equation*}
        d_{c}
        \colon
        \mathrm{R}^{2}\to\mathbb{R}
        \colon
        x\mapsto |x-c|
    \end{equation*}
    which measures the euclidean distance of its argument $x$ from the fixed
    point $c$.
This function is convex thanks to the fact that the euclidean distance is a
    metric and, as such, satisfies the triangle inequality.
However it is not convex strictly since whenever the points $a$ and $b$ lie
    on a common line with the point $c$, the vectors $a-c$ and $b-c$ are
    dependent which boils the term $|\alpha a + (1-\alpha) b - c|$ down into
    an affine function of $\alpha$ as $x\mapsto |x|$ is a norm.

On the other hand, the above mentioned situation which breaks down the
    strict convexity of $d_{c}$ is the \emph{only} case when the function
    exihibits affine behaviour.
All other one dimmensional linear cuts of $d$ are strictly convex.
Indeed, all such cuts are but affine transformations of the function
    \begin{equation}
        x
        \mapsto
        \sqrt{k+x^{2}}
    \end{equation}
    where $k>0$.
Such functions are all clearly strictly convex as one can chech by tedious
    computation, though I am pretty sure it can be done by pure thinking too.
Now, as innocent this might seem it has an interesting consequence.

\begin{proposition}
    Let $c_{1},c_{2}$ and $c_{3}$ be three non colinear points in the plane.
    The function
        \begin{equation*}
            d_{c_{1},c_{2},c_{3}}
            \colon
            \mathbb{R}^{2}\to\mathbb{R}
            \colon
            x\mapsto |x-c_{1}| + |x-c_{2}| + |x - c_{3}|
        \end{equation*}
        is strictly convex.
\end{proposition}
\begin{pf}
    As a sum of three convex functions, this function is clearly convex too.
    In order to conclude that the convexity is also strict let us consider two
        points $a,b$.
    Since the points $c_{i}$ are not colinear, the points $a$ and $b$ can be
        colinear with at most two of the $c_{i}$ or, in other words, at least
        one of the $c_{i}$ is not colinear with $a$ and $b$.
    This means that the term $f(\alpha a + (1-\alpha)b)$ seen as a function
        of $\alpha$ is a sum of three functions of which at most two are
        affine in $\alpha$ while the third one is strictly convex.
    Obviously, sum of an affine function with a strictly convex function is
        a strictly convex function which concludes the proof.
    \qed
\end{pf}

\section{Convex functions related to LED trees}

\begin{proposition}
    The function $f_{1}$ which measures the price of the shortest euclidean
        tree vith leaves $(0,0)$, $(1,0)$ and $x$ and a steiner point $S$
        located optimally (with respect to the total length of the tree) on
        te axis of symmetry of the line segment connecting $(0,0)$ and $x$
        is strictly convex.
\end{proposition}
\begin{pf}
    Consider two points $a$ and $b$ in the plane and the corresponding optimal
        positions of the steiner points $S_{a}$ and $S_{b}$
\end{pf}

\end{document}
